\section{绪论}

\subsection{研究背景}

\par 2024年至2025年间，Solana区块链上的Meme币交易量数次突破历史峰值，单日去中心化交易所交易额一度超越以太坊主网。这场由社区情绪驱动的投机热潮，将Solana的高性能特性展现得淋漓尽致——低费用、高吞吐、秒级最终确认，让普通用户也能以几分钱的手续费参与链上交易。但繁荣的另一面，是自动化程序的泛滥。新代币上线后的前几秒内，常常有数十笔交易同时涌入流动性池，将价格推高数倍后迅速撤出，留下追高的散户承担损失。
\par 这些自动化程序，正是MEV（最大可提取价值）机器人在Solana上的变种。在以太坊和BSC上，研究者已经对狙击机器人、套利机器人的行为模式做了系统梳理，但Solana因架构差异——无公开内存池、并行执行、本地费用市场——始终处于研究的盲区。机器人在Solana上如何运作？它们与在EVM链上有何不同？普通用户能否从交易特征中识别它们？这些问题尚无答案。

\subsection{研究对象}

\par 本文的研究对象是Solana区块链上的MEV机器人账户。所谓MEV机器人，指通过自动化程序执行交易策略、以提取最大可提取价值为主要目的的区块链账户。它们通常表现出高频交易、规律性强、与特定协议深度交互等特征，与人工操作的普通用户形成鲜明对比。
\par 根据策略类型，Solana上的MEV机器人可初步分为以下几类：
\begin{itemize}
    \item 狙击机器人（Sniper Bot）：监控新代币的发行事件或流动性池的创建交易，在新池建立瞬间抢先购买，期望在后续买家涌入时高价卖出。这类机器人与Meme币生态绑定最深，也是前期实验中观察到的活跃主体。
    \item 套利机器人（Arbitrage Bot）：利用不同去中心化交易所之间的价格差，通过低买高卖赚取无风险收益。Solana上的Jupiter聚合器为这类机器人提供了便利的交互接口。
    \item 三明治机器人（Sandwich Bot）：在用户的大额交易前后插入买卖订单，通过影响价格获利。由于Solana无公开内存池，这类机器人的运作机制与以太坊存在差异，可能更多依赖与验证者的合作关系。
    \item 做市机器人（Market Making Bot）：在流动性池中持续提供双边报价，赚取交易费用和流动性激励。这类机器人通常与做市商或量化团队关联，行为模式相对稳定。
\end{itemize}
\par 与以太坊和BSC上的机器人相比，Solana机器人的行为模式存在显著差异。以太坊机器人依赖内存池监控抢跑，Solana机器人则需监控链上事件（如合约部署、流动性池初始化）来捕捉机会；以太坊机器人通过Gas价格竞拍交易顺序，Solana机器人则通过优先费用竞价获取领导者的打包权。这些差异意味着，现有检测方法不能简单移植，必须针对Solana架构做出调整。

\subsection{拟解决的关键问题}

\par 基于上述背景与研究对象的界定，本研究聚焦于以下三个核心问题：
\begin{enumerate}
    \item 如何构建适配Solana架构的机器人行为特征体系？
    \par Solana无公开内存池、并行执行、本地费用市场等特性，使得传统基于内存池监控和Gas费用的检测特征失效。需要重新设计能够捕捉Solana机器人特有行为模式的特征，如优先费用竞价策略、多指令程序调用结构、新代币生命周期参与度等。
    \item 时间序列特征：交易间隔的均值与变异系数、每小时交易分布熵、夜间/周末交易比例、休眠模式特征（如GapBasedSleepiness）等。
    \par 现有研究多将机器人视为同质化对象。但在Solana生态中，狙击、套利、做市等不同类型的机器人，其行为模式、经济动机和生态影响存在显著差异。需要探索能够区分不同机器人亚型的特征与分类方法。
    \item 如何通过消融实验与案例研究增强检测方法的说服力？
    \par 单纯的准确率指标难以证明新特征的必要性，也难以揭示检测结果背后的实际行为含义。需要通过消融实验量化各特征组的贡献，并通过典型案例追踪，将检测结果与链上实际行为相互印证。
\end{enumerate}

\subsection{研究意义}

\par 理论意义：填补Solana生态中自动化账户系统研究的空白。现有MEV研究高度集中于以太坊和BSC等EVM链，对Solana这种非EVM高性能链的关注严重不足。通过揭示Solana机器人的行为特征和运作机制，本研究将拓展MEV研究的跨链视野。同时，针对Solana架构设计的新特征体系，可为其他非EVM链的机器人检测提供方法借鉴。
\par 实践意义：为Solana生态的透明化与健康发展提供技术支撑。通过构建可落地的机器人检测框架，帮助普通用户识别交易对手方性质，辅助监管机构监控市场操纵行为，为DeFi协议设计反机器人机制提供参考依据。

